\documentclass{standalone}
\usepackage{circuitikz}
\usetikzlibrary{calc}
\usetikzlibrary{arrows.meta}
\definecolor{circuitnotered}{HTML}{E5534B}
\definecolor{circuitnoteblue}{HTML}{4C8EDA}
\definecolor{circuitnotemark}{HTML}{FE00FE}
\begin{document}
\begin{circuitikz}[american, line width=0.8pt]
\ctikzset{bipoles/length=1.2cm}
\foreach \x in {0,2,4,6,8} {\foreach \y in {0,-2,-4} {\fill[gray, opacity=0.35] (\x,\y) circle (1.1pt);}}
\node[gray, font=\scriptsize] at (0,0.84) {1};
\node[gray, font=\scriptsize] at (2,0.84) {2};
\node[gray, font=\scriptsize] at (4,0.84) {3};
\node[gray, font=\scriptsize] at (6,0.84) {4};
\node[gray, font=\scriptsize] at (8,0.84) {5};
\node[gray, font=\scriptsize, anchor=east] at (-0.84,0) {a};
\node[gray, font=\scriptsize, anchor=east] at (-0.84,-2) {b};
\node[gray, font=\scriptsize, anchor=east] at (-0.84,-4) {c};
\coordinate (a1) at (0,0);
\coordinate (a2) at (2,0);
\coordinate (a4) at (6,0);
\coordinate (b4) at (6,-2);
\draw (a1) to[R, l_=$R_{1}$, a^=$10\,\mathrm{k}\Omega$] (a2); % line 3
\draw (a4) to[C, l_=$C_{1}$, a^=$100\,\mathrm{n}\mathrm{F}$] (b4); % line 4
\draw[circuitnotered, -{Stealth[length=2.2mm]}] (0,-2) -- (0.598,-0.805); % line 6
\path (-1.365,-2.593) rectangle (1.365,-1.407); % line 7
\node[anchor=west, inner sep=0, circuitnotemark, font=\large] at (0,-2) {X}; % line 7
\draw[circuitnoteblue, -{Stealth[length=2.2mm]}] (6,-4) -- (6,-1.9); % line 8
\path (5.524,-4.593) rectangle (6.476,-3.407); % line 9
\node[anchor=west, inner sep=0, circuitnotemark, font=\large] at (6,-4) {X}; % line 9
\path (10,-7.408) rectangle (19.047,0.593); % line 10
\node[anchor=west, inner sep=0, circuitnotemark, font=\large] at (10,0) {X}; % line 10
\node[anchor=west, inner sep=0, circuitnotemark, font=\large] at (10,-0.487) {X}; % line 10
\node[anchor=west, inner sep=0, circuitnotemark, font=\large] at (10,-0.974) {X}; % line 10
\node[anchor=west, inner sep=0, circuitnotemark, font=\large] at (10,-1.46) {X}; % line 10
\node[anchor=west, inner sep=0, circuitnotemark, font=\large] at (10,-1.947) {X}; % line 10
\node[anchor=west, inner sep=0, circuitnotemark, font=\large] at (10,-2.434) {X}; % line 10
\node[anchor=west, inner sep=0, circuitnotemark, font=\large] at (10,-2.921) {X}; % line 10
\node[anchor=west, inner sep=0, circuitnotemark, font=\large] at (10,-3.408) {X}; % line 10
\node[anchor=west, inner sep=0, circuitnotemark, font=\large] at (10,-3.895) {X}; % line 10
\node[anchor=west, inner sep=0, circuitnotemark, font=\large] at (10,-4.381) {X}; % line 10
\node[anchor=west, inner sep=0, circuitnotemark, font=\large] at (10,-4.868) {X}; % line 10
\node[anchor=west, inner sep=0, circuitnotemark, font=\large] at (10,-5.355) {X}; % line 10
\node[anchor=west, inner sep=0, circuitnotemark, font=\large] at (10,-5.842) {X}; % line 10
\node[anchor=west, inner sep=0, circuitnotemark, font=\large] at (10,-6.329) {X}; % line 10
\node[anchor=west, inner sep=0, circuitnotemark, font=\large] at (10,-6.816) {X}; % line 10
\node[anchor=south west, inner sep=2pt, circuitnotemark, font=\LARGE\bfseries] (circuittitle) at (current bounding box.north west) {X};
\path (circuittitle.south west) rectangle ++(4.298,0.853);
\end{circuitikz}
\end{document}
